\section{신뢰 경계와 재현 방법}
\label{sec:trust-reproduction}

\subsection{무엇을 신뢰하는가}

이 증명 개발의 신뢰 기반(trusted computing base, TCB)은 다음을 포함한다.

\begin{itemize}
  \item EasyCrypt 커널/형 검사기(kernel/type checker)와 가져온 표준 라이브러리;
  \item Why3 및 호출된 SMT 증명기(prover);
  \item Jasmin 구문분석기/컴파일러(parser/compiler)와 \identifier{jasmin2ec}
        번역기;
  \item 해시(hash)로 고정한 Jasmin 소스, 생성 추출물, 재사용 EasyCrypt 산출물;
  \item 소스 해시(source hash) 및 LaTeX 빌드를 수행하는 호스트 도구.
\end{itemize}

\identifier{jasmin2ec}가 소스를 받아들이고 생성된 EasyCrypt가 컴파일된다는
사실은 번역기의 의미보존 증명이 아니다. 별도의 검증된 번역기 정리가 없는 한
번역기 자체는 TCB에 남는다. 마찬가지로 메모리 안전성과 SCT/상수시간
(constant-time)은
EUF-CMA에서 따라오지 않으며 별도의 트랙이다.

\subsection{증명 검증 파이프라인}

저장소 루트에서 다음을 실행한다.

\begin{verbatim}
./haetae-topdown-easycrypt/scripts/verify-all.sh
./haetae-topdown-easycrypt/scripts/check-paper-freeze.sh
python3 haetae-topdown-easycrypt/post-freeze/check-kg-actual-avec-qj-trace.py
\end{verbatim}

첫 명령은 동결 82대상 aggregate를, 둘째는 paper-freeze 범위를 명시적으로,
셋째는 latest post-freeze mismatch의 독립 deterministic witness를 검사한다.

이 스크립트는 다음을 순서대로 검사한다.

\begin{enumerate}
  \item 고정 소스(pinned source)와 읽기 전용 증명 트리(read-only proof tree)의
        SHA-256 변동(drift);
  \item 패커(packer), 서명 코덱, HBZ 준비/적용/전체 래퍼, API 키 복사 보조절차,
        원시 API 호출자, 원시 \(\mu\), 전사 집중 추출(focused extraction)의 재생성과 hash
        일치, Week~7 pack/unpack과 Week~8 focused HBZ procedure body identity 및
        Week~9 실제 rANS 결합 모듈(actual rANS harness)이 재사용하는 생성 대상
        동일성(generated-target identity);
  \item 실제 512단어 HBZ 기호 인증서(symbol certificate)의 결정적 재생성,
        생성기/출력(generator/output) 해시 일치와 생성된 EasyCrypt 인증서 자체의
        커널(kernel) 재검사;
  \item Week~9 rANS 정리 표면에 더해 Week~10 배열/목록, 단어 단계, 내부 진행,
        실제 내부 반복문, 직렬화와 성공 크기 정리 표면 및 Week~11 꼬리 인식
        불변식, 생성 단어 단계, 최종 저장과 실제 인코더 추적 폐쇄, Week~12
        디코더 커서/표/재생 불변식과 실제 추적 정제, Week~13 복사 추적/점별
        입력/W64 크기 어댑터와 실제 핵심 역함수, Week~14 full/focused 내부
        동치, full encoder/decoder, 직접 결합과 production 경계 lift의
        과대주장 검사, Week~15 고정 all-six 예산/실패 원인/production 성공 lift,
        Week~16 실제 M23 matrix/finalizer 스냅샷, 마지막 NTT 표현 rewrite,
        직접 Sign/Verify core control과 \identifier{STOP-VERIFY-MATRIX-CRT}
        동결 표면 검사;
  \item 현재 매니페스트된 82개 작성 EasyCrypt 대상의 \identifier{-no-eco} 새
        컴파일;
  \item \identifier{admit}, \identifier{sorry}, 작성 \identifier{axiom},
        디버그/임시(debug/temporary) 선언 부재와 매니페스트(manifest) 완전성;
  \item 선택된 상류 기준선(upstream baseline)의 새 컴파일(fresh compile);
  \item 기존 주차별(sprint) LaTeX 연구 노트의 참조/인용 오류 없는 빌드;
  \item \sourcepath{manifests/paper-artifacts.md}의 동결 PDF, 82대상, 주장 행렬과
        exact blocker가 변하지 않았는지 확인하는 paper-freeze audit.
\end{enumerate}

{\small\raggedright
추출 해시(hash)는 생성물 변동(drift)을 탐지하지만 번역 의미의 정당성을
대신하지 않는다. 또한 실행 요약은 매 실행 때 덮어쓰므로, 발표용 스냅샷에는
마지막 \texttt{RESULT PASS} 줄까지 포함된 완전한 로그를 보관해야 한다.
76대상 완료 스냅샷의 종료 표식은
\sourcepath{logs/verify-all-before-kg-ntt-mul.log}에 보존된
\texttt{RESULT PASS}\newline
\texttt{authored-targets=76 cache=-no-eco}다. Sign 수정 전 77대상 완료 표식은
\sourcepath{logs/verify-all-before-week16-sign.log}에 보존되고
\sourcepath{WEEK16_SIGN_REPORT.md}는 78번째 대상의 직접 actual-call control과
중단 판정을 설명한다. Verify 수정 전 78/78 완료 로그는
\sourcepath{logs/verify-all-before-week16-verify.log}에 보존되며
해시 \hashsplit{cf8056712327dc8211cf93ae427ac505}{3e8a9d2366747f171392468ac3ff0d75}를
갖는다. \sourcepath{WEEK16_VERIFY_REPORT.md}는 canonical decoded
\((x,v,h,c)\) 경계의 \identifier{STOP-VERIFY-}\allowbreak
\identifier{MATRIX-CRT}와 남은 blocker를
설명한다. paper-freeze 직전 82/82 aggregate 로그는
\sourcepath{logs/verify-all-before-paper-freeze.log}에 보존되고 SHA-256은
\hashsplit{4cd64e5a656be82710bca1410c4d194}{03a3c661d6b91b0319a0ea8f7c91646da}이다.
\newline
기준 커밋 \identifier{e87bcc3}에 추적된
\sourcepath{logs/verify-all-summary.txt}는 paper-freeze audit까지 통과한
terminal \texttt{RESULT PASS}\newline
\texttt{authored-targets=82}\newline
\texttt{cache=-no-eco}와 SHA-256
\hashsplit{08ef9639dc73d56dba42d02999d07897}{d29bc8e60aa100a648e9437fd64387ad}를
유지한다.
동결 논문 PDF 해시는
\hashsplit{be935948028829556951863b44ceb2e6}{c5b2037820991b3a43a4b461b036e53d}다.
\par}

\begin{boundary}[82대상 paper-freeze와 별도 post-freeze 보조 층]
2026-08-10의 Week~14 완료 실행은 Week~13의 67대상 기준선에 다음 다섯 이론을
추가해 72대상을 닫았다.
\begin{itemize}[leftmargin=1.5em,itemsep=0.05em,topsep=0.25em]
  \raggedright
  \item \sourcepath{Mode2HbzInternal}\allowbreak\sourcepath{Boundaries.ec};
  \item \sourcepath{Mode2HbzFullEncodeTrace.ec},
        \sourcepath{Mode2HbzFullDecodeInverse.ec};
  \item \sourcepath{Mode2HbzFullActualInverse.ec},
        \sourcepath{Mode2HbzSignature}\allowbreak\sourcepath{BoundaryLift.ec}.
\end{itemize}
Week~15는 \sourcepath{Mode2RansAllSixBudget.ec}와
\sourcepath{Mode2RansActualSuccessWitness.ec}를 더해 74대상으로, Week~16은
\sourcepath{Mode2KeygenSnapshotAlgebra.ec}와
\sourcepath{Mode2KeygenCoreEquation.ec}를 더해 76대상으로 확장했다.
KG-NTT-MUL 감사가 추가한 paper-freeze 시점 77번째 대상
\sourcepath{Mode2Keygen}\allowbreak\sourcepath{NttMulBridge.ec}는 마지막
\identifier{output_row} 표현
rewrite와 actual two-call harness의 두 active row consequence를 개별 fresh
compile했다. 이 동결 버전은 odd-root convolution이나 security-list
adapter를 가정하지 않으며 판정은 \identifier{STOP-KG-NTT}였다.
\sourcepath{logs/week14-hbz-full-}\allowbreak
\sourcepath{surface-scan.log}는 production 직접 호출,
사전조건, 실패/성공 분기와 frame 결론의 정리 표면을 검사한다. Week~13 독립
검증 기록은 그 아래 실제 rANS 핵심 역함수에 순환 전제가 없음을 별도로
확인한다. \sourcepath{logs/week15-rans-}\allowbreak
\sourcepath{success-surface-scan.log}는 고정 입력
정리가 성공, trace, capacity나 losslessness를 전제하지 않는지 검사하고,
\sourcepath{logs/week16-keygen-}\allowbreak
\sourcepath{core-surface-scan.log}는 실제 두 호출, 비순환
사전조건, NTT boundary가 원하는 곱을 숨기지 않는지와
\identifier{STOP-KG-NTT} 경계를 검사한다.
\sourcepath{logs/week16-sign-core-surface-scan.log}는 실제 세 helper 호출,
비순환 precondition과 stop 경계를 검사한다. 실행 중 생성되는
\sourcepath{logs/verify-all-summary.txt}의 중간 PASS 줄은 완료 증거가 아니며,
terminal RESULT와 고정 해시가 함께 있을 때만 현재 aggregate 증거로 쓴다.
Verify 수정 전 78/78 완료 로그는
\sourcepath{logs/verify-all-before-week16-verify.log}에 보존되며,
paper-freeze 직전 82/82 aggregate는
\sourcepath{logs/verify-all-before-paper-freeze.log}에 보존된다.

post-freeze의 \sourcepath{NTTFullSpectralAction.ec},
\sourcepath{NTTMode2RowSpecializations.ec},\newline
\sourcepath{RqHAETAEBridge.ec}는 외부 NTT 트리에서 fresh compile됐고, 강화된
\newline\sourcepath{Mode2KeygenNttMulBridge.ec}는 현재 82대상 aggregate 안에서 다시
컴파일된다. 세 \sourcepath{post-freeze/KgActualAvecQj*.ec} 대상은 동결
매니페스트에 넣지 않고 각각 \identifier{-no-eco}로 별도 컴파일했다. 따라서
82라는 수는 post-freeze 보조 정리까지 모두 센 수가 아니며, 동결 proof-target
surface가 유지됐다는 뜻이다.
\end{boundary}

\subsection{이 문서 빌드}

이 안내서는 기존 sprint 노트와 독립된 XeLaTeX 문서다. 저장소 루트에서

\begin{verbatim}
sh haetae-topdown-easycrypt/theory-guide/build.sh
\end{verbatim}

을 실행하면 \sourcepath{haetae-topdown-easycrypt/theory-guide/main.pdf}가
생성된다. 스크립트는 중단 오류와 정의되지 않은 참조/인용(undefined
reference/citation)을 실패로 처리한다.

문서 작성 시 관찰된 도구 경계는 Jasmin/\identifier{jasmin2ec} 2026.03.0,
EasyCrypt OPAM switch \identifier{easycrypt-5.2}, Why3 1.8.0이며, 구성된
prover에는 Alt-Ergo 2.6.0, CVC4 1.8.0, CVC5 1.2.1, Z3 4.12.6이 포함된다.
정확한 재현 결과는 버전 문자열보다 실행 로그를 우선한다.

\begin{takeaway}{완료 판정 기준}
이 안내서가 컴파일되는 것과 EasyCrypt 정리가 컴파일되는 것은 서로 다른
증거다. 문서 배포 전에는 \identifier{build.sh}와 전체
\identifier{verify-all.sh}를 모두 성공시켜야 한다. 현재 동결 증명 기준선은
82/82 aggregate와 paper-freeze audit의 terminal PASS다. post-freeze 주장은 세
actual-\identifier{avec}/paper-\identifier{qj} 대상의 개별 fresh compile, Rq/NTT 보조
대상과 deterministic trace까지 별도로 확인해야 한다. 덮어쓰기 중인 summary의
중간 PASS 줄은 완료 증거가 아니다. 이 안내서의 별도 빌드도 함께 확인한다.
\end{takeaway}
